\pdfmapfile{+cm.map}
\pdfmapfile{+cmextra.map}
\pdfmapfile{+symbols.map}
\documentclass[11 pt]{article}
\usepackage[margin=0.86 in]{geometry}
\usepackage{amsmath,amssymb,amsthm,booktabs,array,hyperref}
\hypersetup{colorlinks=true,urlcolor=blue,linkcolor=blue,pdftitle={Twelve positive unknowns and materialized coefficients for ZFC consistency},pdfauthor={Prepared for Vladimir Reshetnikov}}
\setlength{\emergencystretch}{3 em}
\newcommand{\Pos}{\mathbb Z_{>0}}
\newcommand{\Con}{\operatorname{Con}}
\newtheorem{theorem}{Theorem}
\newtheorem{lemma}{Lemma}
\title{Twelve positive unknowns and materialized coefficients for ZFC consistency}
\author{Research report prepared for Vladimir Reshetnikov}
\date{22 October 2026}
\renewcommand{\labelitemi}{$\bullet $}
\begin{document}
\maketitle
\begin{abstract}
We improve the latest thirteen-positive-variable construction to twelve positive variables by combining its remaining main Pell root with its existing three-root witness. A separation bound and a rational auxiliary-root integrality lemma exclude the spurious solutions that a resultant alone would permit. The resulting polynomial has exact degree 1145 and a 123-operation circuit. Separately, a sharper root weight gives a 27-unrestricted-integer construction; low-degree lifts reach 54 integer variables at degree 4,45 at degree 6 and 39 at degree 8. A new 4094-tile binary-copy compiler lowers the base to $ 2^{5^{11}}$. All four fixed coefficient integers, totaling 25,421,252 bytes, are supplied as actual binary files, not just finite recipes. These are mathematical reductions with reproducible checks, not newly kernel-checked composite proofs or global minimality results.
\end{abstract}
\section{Results and comparison with the recent commits}
The inspected source snapshot is
\texttt{f0ae4e0991ca44a7bb4e5671d42e33d775e48035}.
Its wave-2 collection contains packages 09--15, superseding the old 01--08 directories. We use 09's root amalgamation, 12's low-degree lifts, 14's signed-domain and square-scale arguments, 15's low-degree parents, and the earlier 11 T4 compiler as explicit inputs. Historical references remain pinned; removed directories are not assumed to exist at the new snapshot.

\begin{center}\small
\begin{tabular}{lrrrr}
\toprule
New construction&Unknowns/domain&Equations&Degree&Operations\\
\midrule
Main root product&12 positive&1&1145&123\\
Main root system&12 positive&3&1136 maximum&119\\
Smaller-weight product&13 positive&1&515&84\\
Signed small-weight product&27 integers&1&1004&136\\
Signed quadratic lift&54 integers&31&2&not optimized\\
Its sum of squares&54 integers&1&4&not optimized\\
Sextic alternative&45 integers&1&6&not optimized\\
Octic alternative&39 integers&1&8&not optimized\\
\bottomrule
\end{tabular}
\end{center}
The 12-variable result improves the latest 13-variable frontier in the supplied collection, but trades a larger degree and circuit for the saved coordinate. The new 13-variable circuit preserves 84 operations while lowering degree 571 to 515; the previous degree 515 variant cost 85. The 27-integer construction improves the 29-variable count but not its degree 188 or 121-operation cost. The 54-integer quartic improves the reported 60-integer quartic without changing degree. These are different objectives, not one uniformly dominating result.

The inherited 67-operation product remains the best operation count in this comparison. We have not proved a global minimum, and no comparison with abstract universal-polynomial minima is claimed. An operation here is a binary addition, subtraction or multiplication, including arithmetic by fixed numerals; shared reuse and comparisons are free. Coefficient bit length is not included in circuit counts, and is reported separately.

All the constructions have the same intended theorem:
\begin{equation}\label{eq:goal}
\mathrm{ZFC}\vdash
\bigl(\text{no solution in the explicitly stated domain}\bigr)
\longleftrightarrow\Con(\mathrm{ZFC}).
\end{equation}
No consistency assumption is needed to prove the equivalence. It does not prove either side.

\section{Fixed integers and the common arithmetic interface}
The main instance uses the complete four-tuple in
\texttt{coefficients/COEFFICIENTS.json}, with actual values in
\texttt{coefficients/expanded/}. A file is interpreted as an unsigned, minimal-width, big-endian integer. Its unused leading bits are zero. The exact file hashes and populations are in \texttt{expanded/MANIFEST.json}.
\begin{center}\small
\begin{tabular}{lrrr}
\toprule
File/constant&Exact bit length&Bytes&Set bits\\
\midrule
\texttt{beta.bin}: $\beta $&48828125&6103516&48828125\\
\texttt{K.bin}: $ K $&56885626&7110704&133545\\
\texttt{MC.bin}: $ M_C $&48828125&6103516&48823982\\
\texttt{MF\_source.bin}: $ M_F^*$&48828126&6103516&4145\\
\bottomrule
\end{tabular}
\end{center}
Here $\beta=B-1 $, $ K=D_C+BD_R $, and $ M_F^*=M_F^{\rm nat}+B-1 $ is the \emph{shifted} high mask. The native high mask has 4143 set bits. Do not replace the source port by the native mask. These four numbers are fixed, not additional unknowns. Their finite generating formulas and complete tile list are supplied as well as their expanded binary values. The full polynomial is represented by a factored circuit, not an expanded monomial list.

The literal 215-state binary K215 table has SHA-256
\begin{center}\small
\texttt{fdbcef0825886035c3091c8537a69535bc96f9e265f0ae45863543d40ae70090}.
\end{center}
It starts at state 0 on blank symbol 0 and halts at state 215. Its real first action is $(0,0)\mapsto(1,+1,93)$. Only \emph{after} entering halt do stationary identity transitions provide tableau padding. None of the 430 original transitions is changed.

For the twelve positive unknowns
\[
(J,F,\alpha,z,h,i,T,s,w,\eta,\zeta,t)\in\Pos^{12},
\]
define substituted polynomial abbreviations
\begin{align*}
q&=\beta J+1,&C_w&=q-F-\alpha,&X&=wq,&Y&=sq^3,\\
E&=XY,&k&=\eta+\zeta,&c&=kY+\eta,&a&=E+Y,\\
H&=4 a+3,&\Delta&=a^2+H,&&&
\end{align*}
and
\begin{align*}
R&=((q-1)(q-F)+\alpha)(q^2-1)+(M_C+qM_F^*)J,\\
N_k&=k-hE-R,\\
N_t&=(K+w)C_w+q-F-z(q-1).
\end{align*}
These are abbreviations, not additional variables or equations. All exponents are fixed.

\section{The explicit twelve-variable polynomial}
Introduce further polynomial abbreviations
\begin{align*}
L_0&=EkY,& A&=1+L_0(L_0+k),& I&=ic^2,\\
D&=1+\Delta I^2,& Q&=\Delta(D-1)=(\Delta I)^2,&d&=Q-1,\\
v&=cT,&z_0&=c+RD,&P&=Dv^2+z_0^2,\\
e&=QP-1,&B_a&=2 Qvz_0,&M&=2(kI^2+T^2).
\end{align*}
Set
\[
S=2 M(L_0+k),\qquad \Lambda=HS+1,\qquad
b_0=X+ac,\quad G=1+\Delta c^2,\quad W=Ht+\Lambda b_0.
\]
The following expressions are an explicit even/odd norm calculation; they use only polynomial operations, not square roots or a black-box resultant:
\begin{align*}
Z_0&=W^2+\Lambda^2 G,& Z_1&=-2 W\Lambda,& n&=H^2 M^2 A,\\
C_0&=d(Z_0+n)-H^2 e,& C_1&=dZ_1,&&\\
q_0&=C_0^2+GC_1^2,& q_1&=2 C_0 C_1,&&\\
j_0&=H^4 DB_a^2,&r_0&=4 d^2 n,&&\\
U_0&=q_0+j_0-r_0 Z_0,&U_1&=q_1-r_0 Z_1,&&\\
A_e&=U_0^2+GU_1^2-4 j_0 q_0,& A_o&=2 U_0 U_1-4 j_0 q_1.&&
\end{align*}
Finally define
\begin{equation}\label{eq:res}
\mathcal R=A_e^2-GA_o^2,
\end{equation}
and the requested single polynomial
\begin{equation}\label{eq:p12}
\boxed{P_{12}=N_kN_t(1+\mathcal R)-1.}
\end{equation}
The alternative three-equation system is
\[
N_k=1,\qquad N_t=1,\qquad\mathcal R=0.
\]
Its sum of squared residuals gives another single equation of degree 2272. The attached JSON files list every binary gate; the circuit shares repeated subexpressions and costs 123 operations for \eqref{eq:p12}, not the cost of naively expanding these displays.

\begin{theorem}\label{thm:12}
Existence of a positive integer zero of $ P_{12}$ is equivalent to blank K215 halting, hence to a finite ZFC refutation. The same holds for the three-equation system and its sum of squares. The reductions are formalizable in ZFC.
\end{theorem}
\subsection{Bounds before recovering any root}
At a zero of \eqref{eq:p12}, the three factors are integer units. Since $\Delta\equiv 0,3\pmod 4 $, $ G\equiv 0,1\pmod 4 $, so $\mathcal R $ cannot be 2 modulo 4. Therefore $ 1+\mathcal R $ cannot be $-1 $, and
\[
\mathcal R=0,\qquad N_k=N_t=\varepsilon\in\{1,-1\}.
\]
Weak transport and the authentic mask bounds give
\[
C_w\ge 0,\quad q\ge 16,\quad 48<R<q^4-2,\quad E\ge q^4>R+2,
\quad X\ge 16,\quad Y\ge 4096.
\]
Because $ h>0 $, $ k=R+\varepsilon+hE>E $ and $ k>R $. The possible $ C_w=0 $ branch is not prematurely removed.

Let $ A_p=a+2 $ and $ M_0=kI^2+T^2 $. The elementary estimates
\[
c<2 kY,\quad A_p<2 E,\quad L_0>A_pc/4,\quad L_0>4 A_p^2
\]
imply
\[
A>A_p^2 c^2/16,\quad A>16 A_p^4,\quad D<A_p^2 I^2,
\quad z_0<2 kA_p^2 I^2.
\]
Together with $ M_0^2\ge 4 kI^2 T^2 $ and $ M_0\ge kI^2 $, these give
$ AM_0^2>4 Dv^2 $ and $ AM_0^2>4 z_0^2 $. Consequently
\begin{equation}\label{eq:domination}
dAM^2>4 QP.
\end{equation}
This bound uses no main Pell root or auxiliary integrality premise.

\subsection{The inner root polynomial and complex separation}
For a scalar $ x $, define
\begin{align*}
C(x)&=d(x^2+M^2 A)-e,\\
U(x)&=C(x)^2+DB_a^2-4 d^2 x^2 M^2 A,\\
\mathcal F(x)&=U(x)^2-D(2 C(x)B_a)^2.
\end{align*}
For real $ x $, \eqref{eq:domination} implies
\begin{equation}\label{eq:cpositive}
C(x)>3 QP+1>B_a\sqrt D>0.
\end{equation}
Factoring twice over the complex numbers gives the complete root description
\[
x=M\tau+y,\qquad \tau^2=A,\quad f^2=D,\quad
 d y^2=Q(vf-z_0)^2-1.
\]
In particular, for every branch, even if $ y $ is imaginary,
\[
|y|^2\le(2 QP+1)/d<4 QP/d<AM^2.
\]
Since $ A<(L_0+k)^2 $, \emph{every} complex root of $\mathcal F $ obeys
\begin{equation}\label{eq:rootbound}
|x|<2 M\sqrt A<S.
\end{equation}
It would be incorrect to assume every conjugate auxiliary root is real.

Let $\mathcal F_H(z)=H^8\mathcal F(z/H)$. Direct expansion of the displayed even/odd expressions gives
\[
\mathcal F_H(W-\Lambda b)=A_e+bA_o\quad\text{when }b^2=G.
\]
Thus \eqref{eq:res} is the product of this expression for $ b=\pm\sqrt G $. If $ G $ were nonsquare, $ A_e^2=GA_o^2 $ with integral $ A_e,A_o $ would force both to vanish. Both
\[
x_\pm=(W\pm\Lambda\sqrt G)/H
\]
would then be roots of $\mathcal F $, contrary to
\[
|x_+-x_-|=2\Lambda\sqrt G/H>2 S.
\]
Hence $ G $ is a square. Choose an integer $ b $ such that
$ x=(W-\Lambda b)/H $ is a root. It is positive $ b $: otherwise
$ W=\Lambda b+Hx<-\Lambda\sqrt G+HS<0 $, while $ W>0 $.

\subsection{Integrality, not merely a resultant zero}
The rational $ x $ is nonzero, since \eqref{eq:cpositive} gives
$\mathcal F(0)=(C(0)^2-DB_a^2)^2>0 $. Put
\[
f=-\frac{U(x)}{2 C(x)B_a},\qquad
\tau=\frac{C(x)+B_af}{2 dxM},\qquad y=x-M\tau.
\]
All denominators are established nonzero. We obtain $ f^2=D $ and $\tau^2=A $, so the rational numbers $ f,\tau $ are integers. The remaining equation is
\[
Q(vf-z_0)^2-(Q-1)y^2=1.
\]
The following lemma is the essential missing step in a naive root amalgamation.

\begin{lemma}[Rational auxiliary-root integrality]\label{lem:rational}
If $ J_0,V\in\mathbb Z $, $ J_0\ge 2 $, $ y\in\mathbb Q $, and
$(J_0 V)^2-(J_0^2-1)y^2=1 $, then $ y\in\mathbb Z $.
\end{lemma}
\begin{proof}
Write $ J_0^2-1=d_0 s_0^2 $ with $ d_0>1 $ squarefree. The rational number $ z=s_0 y $ satisfies $ d_0 z^2\in\mathbb Z $; a reduced denominator would have its square divide the squarefree $ d_0 $, so $ z $ is integral. The positive integral Pell group gives indices $ m,n\ge 1 $ with
\[
J_0=x_m,\quad s_0=z_m,\quad |J_0 V|=x_n,\quad |z|=z_n,
\]
where $ x_r+z_r\sqrt{d_0}$ are powers of the fundamental positive integral unit. Modulo $ x_m $, this unit's $2m$-th power is $-1 $. Therefore
$ x_{r+2 m}\equiv-x_r $ and $ x_{2 m-r}\equiv-x_r $.
Reducing and reflecting the index, $ x_m\mid x_n $ and
$ 1\le x_r<x_m $ for $ r<m $ force $ n $ to be an odd multiple of $ m $.
The multiplication formula gives $ z_m\mid z_n $, hence $ s_0\mid z $ and $ y $ is integral. This uses integer Pell solutions, not half-integral units of a larger ring.
\end{proof}
Apply the lemma with $ J_0=\Delta I $ and $ V=vf-z_0 $. It follows that $ x=M\tau+y $ is integral. Since
$ W=Ht+\Lambda b_0 $ and $ W-\Lambda b=Hx $, and since $\Lambda=HS+1 $, we recover
\[
\gamma=(b-b_0)/H\in\mathbb Z.
\]
The main Pell equation and its congruence are now restored exactly.

In fact $\gamma>0 $ follows directly. The first norm classifies
$ k=2\psi_{2 XY^2+1}(n)$. Since $ k>R>48 $, $ n\ge 2 $ and $ c>k>X $.
Also $ b^2-\Delta c^2=1 $, $ b>0 $, and $\Delta>(a+1)^2 $ give
$ b-ac>c>X $. Hence $\gamma H=b-ac-X>0 $.

The recovered $ f,\tau,y $ are nonzero but may initially be signed. Package 09's explicit normalized signed-$ f,T $ argument applies: first/main Pell indices, signed auxiliary congruences, half-interval step-down and no-wrap recover the packed index, dyadic scale and native tableau. It does not require assuming the reconstructed auxiliary quotient positive. Its exact proof notes are included. This proves soundness without silently assuming positivity of a recovered root.

Conversely, start with a canonical positive normalized blank witness and set
\[
t=\Lambda\gamma+M\tau+y>0.
\]
Then $ W=\Lambda b+H(M\tau+y)$ and one factor $\mathcal F_H(W-\Lambda b)$ vanishes. The two transport/index factors are 1. This proves completeness and Theorem \ref{thm:12}.
\subsection{Exact degree, including the leading cancellation}
The large degree is a genuine cost of saving a coordinate, but it is smaller than the naive degree of the displayed norm expression. The base degrees are
\[
\deg c=5,\quad\deg a=6,\quad\deg\Delta=12,\quad
\deg M=23,\quad\deg S=34,\quad\deg\Lambda=40,\quad\deg W=51.
\]
The important cancellation is
\[
b_0^2-G=X^2+2 Xac-Hc^2-1,
\]
which has exact degree 16, not 22. Hence
\[
K_p=W^2-\Lambda^2 G
=\Lambda^2(b_0^2-G)+2 Ht\Lambda b_0+H^2 t^2
\]
has exact degree 96.
Write $\mathcal F_H(z)=\sum_{j=0,2,4,6,8} a_jz^j $. Then
$ a_8=d^4 $ and $\deg a_j\le 536-44 j $. Pairing conjugate cross terms in the norm yields the degree bound
\[
1072+7\max(i,j)+\min(i,j)
\]
for the $(i,j)$ term. The unique maximum is at $ i=j=8 $, with nonzero leading form $ d_{\rm top}^8(K_p)_{\rm top}^8 $. Therefore
\[
\deg\mathcal R=1136,\qquad\deg P_{12}=1136+7+2=1145.
\]
The sum-of-squares degree is 2272. Uniform leading forms prove these statements for the actual coefficients; finite-field line diagnostics are supplementary checks, not the proof of a uniform degree claim.

\section{A cheaper thirteen-variable weight and 27 integer unknowns}
A separate improvement retains the main-root parameter $\gamma $ and amalgamates only $ f,\tau,y $, as package 09 did. Use the same inner polynomial $\mathcal F(t)$ but replace its weight by
\[
\boxed{M=I^2+T^2.}
\]
Its positive-domain product is
\[
\bigl((X+ac+\gamma H)^2-\Delta c^2\bigr)
N_kN_t(1+\mathcal F(t))-1.
\]
This costs 84 operations and has exact degree 515, improving the previous 84-operation degree 571 construction without adding a variable. Unlike the twelve-variable construction, this proof can use the main norm \emph{before} establishing the weight bound. The smaller weight must not be transplanted into the twelve-variable proof without that missing premise.

Here are the new bounds. Let $ A_p=a+2 $ and $\lambda=(16/15)^4 $. The main norm gives $ c=\psi_{A_p}(p)$, with the absolute main root if necessary. The index $ p=1 $ would give $ c=1 $. At $ p=2 $ its congruence gives either $ X\equiv 4\pmod H $ or $ 4 X\equiv-1\pmod H $, both impossible because $ X\ge 16 $ and $ 0<4 X<H-1 $. Thus $ p\ge 3 $, and
\[
c\ge 4 A_p^2-1,\qquad k>\tfrac 72 X^2 Y>E.
\]
With the weak rank bounds, one obtains
\[
\frac{Dc^2}{I^2}<\lambda A,\qquad
\frac{z_0^2}{I^4}<\frac{4\lambda}{49}A.
\]
Consequently
\[
\frac{Dv^2+z_0^2}{(I^2+T^2)^2}
<\lambda\left(\tfrac 14+\tfrac 4{49}\right)A<\tfrac 9{20}A.
\]
Since $ Q\ge 10 $ and $ d\ge 9 Q/10 $, this proves
$ dAM^2>2 Q(Dv^2+z_0^2)$ and hence
$ C(t)>|B_a|\sqrt D $ for every real $ t $.

\subsection{The signed-domain extension}
Only the seven ports
\[
J,F,\alpha,z,s,\eta,\zeta
\]
need remain positive. Put $ i=2 u+1 $ and $ w=v_w^2 $, with $ u,v_w\in\mathbb Z $; allow $ h,T,\gamma,t $ to be arbitrary integers.
The zero-square boundary is excluded first: if $ w=0 $, then $ X=0 $ and the main norm modulo $ H $ gives $ 0\equiv 1\pmod H $, impossible. Thus $ w>0 $ before growth estimates are used.

Two other degeneracies need proofs. If $ t=0 $, the resultant is
$(C(0)^2-DB_a^2)^2>0 $ by strict domination. If $ T=0 $, then $ B_a=0 $; since $ t\ne 0 $, the resultant still recovers an integral first root and an integer auxiliary root, giving
\[
Qz_0^2-(Q-1)y^2=1.
\]
The outer bounds imply $ 1<z_0<Q $. But any such Pell solution with $ z_0>1 $ requires $ z_0\ge 2 Q-1 $, a contradiction. Thus both denominators are nonzero. The original root reconstruction and normalized signed proof now apply.

For completeness, canonical padding uses odd powers of five, so $ dN $ is odd and the tuned rank $ R\equiv 3\pmod 4 $ is odd. The scale
$ w=2^{R-dN}$ is therefore a positive square. The strong auxiliary index can be chosen $ cR $, yielding the odd integer
\[
i=\psi_{a+2}(cR)/c^2.
\]
The old positive auxiliary construction restores all compressed quotient congruences, and $ t=M\tau+y $ is a positive integer. This is a fresh canonical witness argument, not an assertion that every existing $ i $ is odd or every positive $ w $ is square.

Replace each of the seven positive ports by an independent copy of
\[
\operatorname{Pos}(r_1,r_2,r_3)=1+r_1^2+r_2^2+r_3(r_3+1).
\]
It parametrizes all positive integers, by the three-square theorem applied to $ 4(n-1)+1 $. The total is
$ 7\cdot 3+1+1+4=27 $ unrestricted integer variables. The explicit product costs 136 operations and has exact degree 1004; its four-equation system costs 130 operations with maximum degree 952. A single SOS costs 141 operations and has degree 1904. All substitutions, gate counts and nonzero leading coefficients are supplied.

\section{Smaller low-degree equations over unrestricted integers}
The preceding constructions minimize witness count at the expense of degree. A different combination improves the low-degree choices.

Start with package 12's 32-positive-coordinate, 23-quadratic lift. Replace the strong multiplier $ i $ affinely by $ 2 r+1 $. Add one integer $ u $ and the equation $ w-u^2=0 $. For each of precisely seven guarded coordinates $ J,F,\alpha,z,s,\eta,\zeta $, add three integer parameters and a quadratic equation
\[
x-1-r_{x, 1}^2-r_{x, 2}^2-r_{x, 3}(r_{x, 3}+1)=0.
\]
All other old coordinates are unrestricted. There are
\[
32+1+7\cdot 3=54
\]
integer variables and $ 23+1+7=31 $ quadratic residuals. Their sum of squares is a quartic $ Q_{54}$. The coefficient of $\tau^4 $ is exactly 1, so its degree is exactly 4. Every residual and the expanded sparse polynomial are in \texttt{domains/integer54\_quartic.json}.

This lift requires more than dropping positivity annotations. The proof restores the main arithmetic, derives $ w>0 $ from the main norm, applies the ordinary strong-rank theorem, then obtains $ f^2\equiv 1\pmod c $. Only then is the merged quotient $ T=(o+Rf)/c $ known integral. The transport coordinate in this particular parent is shifted by one; the proof retains that convention. The full projection and positive canonical lift are included.

Applying the same guard mechanism to package 15's lower-degree parent systems, while substituting guards directly when the residual degree permits, yields:
\begin{center}\small
\begin{tabular}{lrrr}
\toprule
File&Integer variables&Residuals&SOS degree\\
\midrule
\texttt{integer54\_quartic.json}&54&31&4\\
\texttt{integer45\_degree6.json}&45&22&6\\
\texttt{integer39\_degree8.json}&39&16&8\\
\bottomrule
\end{tabular}
\end{center}
The last two exact-degree certificates are the positive coefficients
$\beta^2 $ of $ J^2 L^2 s^2 $, and $\beta^6 $ of $ J^6 s^2 $, respectively. Their literal variable lists and sparse coefficients are supplied, not merely the existence of some polynomial lift. These are unrestricted-integer forms. They must not be conflated with 32-positive-variable quartics or with the 27-integer high-degree equation.

\section{Actual coefficient integers from 4094 local tiles}
This improvement is independent of all the arithmetic reductions. It replaces the complete compiler tuple, not individual ports, and reduces the bit length of $\beta $ by a factor of 78125 compared with the preceding T4 tuple. It also makes materialization practical.

\subsection{Safe filtering of the local language}
The old T4 cell codes are retained: 439 symbols, including two headless tape bits, headed bits for all 216 states including padded halt, three initialization phases and two boundaries. A tile contains left, center, right, next-time-center. Start is $(434,435,436,1)$; the forbidden selector is $(438,0,0,438)$, locally valid under the old predicate but globally impossible.

The new language is a subset of old T4 tiles, plus that explicit forbidden tile. It imposes compatible adjacent phase/boundary pairs, blank boundary collars, no immediate initial-to-final boundary, and a simple predecessor condition: an ordinary headed cell must have some original incoming transition whose written source bit matches the observed headless source cell whenever that cell lies in the triple. This condition deliberately uses the 430 original transitions, not halt padding transitions.

Soundness follows because every new global tableau is an old one. For completeness, enlarge each finite genuine run by at least three permanently blank cells between the head's visited region and either boundary. Every genuine ordinary head has its original incoming transition. During halt padding, the source cell of the original halting transition remains unchanged and headless. The collars protect the helical seams. Width and height remain independently paddable to sufficiently large powers of five.

The ordered alphabet is Start, forbidden, then the remaining valid tuples lexicographically. It has 4094 entries, and is materialized as\texttt{tiles4094.json}. Independent enumeration and a different center-class count agree. The complete predicate is executable and finite; it does not wait for K215 to halt.

\subsection{Binary copies and occupancy}
Replace each 439-way one-hot symbol copy by the eleven-bit code
\[
\chi(s)=(\operatorname{bit}_0(s),\ldots,\operatorname{bit}_8(s), 1,0).
\]
The nine payload bits are injective; the constant 1 prevents an occupied blank symbol from masquerading as an empty cell; the zero guard preserves convenient layout arithmetic. Use $ A_c=4096 $ and
\[
\mu=A_c-2+\sum_{j=1}^{48}A_c^j.
\]
Before occupancy is proved, each clause digit has at most 4095 contributions, below $ A_c $, so no base-$ A_c $ carry is possible. The selector-count mask forces zero or one selector. Parity then forces its matching copies and anchors. The occupancy bit propagates through the spatial overlap around the cyclic word. Thus binary coding preserves the old exact-symbol and occupancy theorem, not merely bitwise necessary conditions.

\subsection{Exact finite recipe and the new zero padding}
Set $\rho=2^{625}$ and $ B=\rho^{78125}=2^{5^{11}}$. The native support contains selectors $ 0,\ldots, 4093 $, four eleven-bit copy blocks starting at
$ 4182,4171,4160,4138 $, one dummy at 4226, and anchors
$ 4260,4272,4282,8565 $. Its size is 4143.
For tile $ i $ with entries $ w_i[0],\ldots,w_i[3]$, let
\[
v_i=1+\sum_{j=0}^3\sum_{r=0}^{10}\chi(w_i[j])_r A_c^{1+11 j+r}
+\sum_{r=1}^4 A_c^{44+r}.
\]
A copy bit has the corresponding monomial coefficient, the four anchors have $ A_c^{45},\ldots,A_c^{48}$, and the dummy coefficient is zero. If $\mathcal E $ is the support, define
\begin{align*}
D_C={}&\rho^{33}+\rho^{12903}+\rho^{22}+\rho^{4293}
 +\sum_{e\in\mathcal E}v_e(\rho^{30034-e}+\rho^{47168-e}),\\
D_R={}&\rho^{12892},\\
M_F^0={}&\mu(\rho^{30034}+\rho^{47168})+\sum_{r=0}^{10}\rho^{4171+r}\\
&+\sum_{r=0}^{10}(\rho^{12892+4182+r}+\rho^{12892+4171+r})
 +\rho^{4282}+\rho^{8565}+\sum_{z=55734}^{59722}\rho^z,\\
M_C={}&B-1-\sum_{e\in\mathcal E,\ e\ne 1}\rho^e-2\rho^{4226},
\qquad M_F^{\rm nat}=M_F^0+4.
\end{align*}
Finally $\beta=B-1 $, $ K=D_C+BD_R $, and $ M_F^*=M_F^{\rm nat}+B-1 $.

The important new idea is the final 3989 mask bits. They balance populations without making $\mu $ enormous. Every possible aligned field contribution has radix exponent at most $ 47168+8565=55733 $. The conservative carry bound prevents overflow into 55734. Thus the added tests are identically zero for every intended field, including every allowed dummy adjustment. They do not constrain completeness. Soundness may ignore these extra restrictions.

Actual integer checks give
\[
\operatorname{pc}(M_C)=48823982,\qquad
\operatorname{pc}(M_F^{\rm nat})=4143,
\qquad \operatorname{pc}(M_C)+\operatorname{pc}(M_F^{\rm nat})=48828125.
\]
Also $ M_C\equiv 2\pmod 4 $ and $ M_F^{\rm nat}\equiv 4\pmod 8 $. The mass threshold has 604 bits and $\mu $ has 577 bits, both safely below the 625-bit radix.

The two clause bands and two tested anchors still exclude every non-whole-cell rotation, including the optional upper-dummy spill. Independent exact checks cover 435015 basis-convolution cases, all relevant support rotations, and the actual finite mask data. The compact anchor geometry is proved explicitly in the coefficient proof.

The five-adic unit is now
\[
2 D_C-D_R\equiv 1\pmod 5,
\]
not the old value 3. Only nonzero residue is needed for actual-index tuning. Because width, height and $ d $ are powers of five, the old subset-adjustment proof applies. The extra zero fields remain zero after each dummy change. Canonical $ R\equiv 3\pmod 4 $ and odd $ dN $ survive, so the odd auxiliary and square-scale completeness arguments remain valid. This establishes the whole compiler interface needed by every arithmetic form in this report.
\section{Proof dependencies, reproduction, and recommendation}
The proof chain is:
\[
\begin{gathered}
\text{a solution of the chosen polynomial}\ \Longleftrightarrow\
\text{the normalized Pell/rank system},\\
\Longleftrightarrow\ \text{a marked compact T4 tableau}
\ \Longleftrightarrow\ \text{a blank K215 halt}
\ \Longleftrightarrow\ \text{a ZFC refutation}.
\end{gathered}
\]
The new root-elimination, signed-domain and compiler-adapter arguments are provided here and in the mathematical proof files. The underlying shifted-population compiler, Pell decoding and K215 logical calibration are substantive imported theorems, not assumptions that K215 fails to halt. The supplied arithmetic audits do not by themselves certify those unbounded implications. Neither a new Lean/Rocq end-to-end build nor a formal ZFC proof-code is claimed.

The archive includes literal arithmetic circuits, expanded sparse low-degree equations, their generators, independent algebraic checks, exact tile enumeration, the four materialized coefficient integers, and pinned source evidence. The root reproduction driver verifies the shipped data before rebuilding in a temporary copy. The README gives the supported command, dependency list and exact scope. Regeneration uses no K215 halting search and does not require constructing an actual refutation witness.

Independent checks exercise the displayed identities, exact degree certificates and finite compiler obligations. They do not replace the proofs or establish $\Con(\mathrm{ZFC})$. The proofs explicitly handle zero denominators, signed witnesses and irrational outer roots.

For incorporation, use the twelve-positive-variable form when witness count is the main objective, the 54-integer quartic when low degree and an unrestricted domain matter, and the thirteen-positive-variable form when a smaller circuit matters. Keep the complete compact coefficient tuple as a shared backend. Preserve the older low-gate-count and low-degree variants rather than replacing the whole catalogue with a single ``winner.'' Further progress may come from combining the outer-root amalgamation with the signed-domain argument, but the present twelve-variable proof depends on positivity before root recovery; simply dropping that guard is not justified.

\subsection*{Pinned source locations}
All repository comparisons refer to Vladimir Reshetnikov's ProveIt repository. The source evidence in the archive retains exact paths and hashes.
\begin{itemize}
\item \href{https://github.com/VladimirReshetnikov/ProveIt/tree/f0ae4e0991ca44a7bb4e5671d42e33d775e48035/SetTheory/Cardinals/docs/reports/hilbert-tenth-problem/con-zfc-diophantine-forms}{Diophantine comparison collection at the reviewed commit}, especially packages 09,11,12,14 and 15.
\item \href{https://github.com/VladimirReshetnikov/ProveIt/tree/f0ae4e0991ca44a7bb4e5671d42e33d775e48035/Computability/HilbertTenthProblem}{Hilbert's tenth problem compiler sources}.
\item \href{https://github.com/VladimirReshetnikov/ProveIt/tree/f0ae4e0991ca44a7bb4e5671d42e33d775e48035/Computability/BusyBeaver/Research/ZFC/k225-formal/verification/k215}{K215 machine and verification sources}.
\end{itemize}
\end{document}
